% Cleo bench — shared template for derivation documents.
% Replace <<TITLE>>, <<QID>>, <<N>> and the body. Keep the preamble as-is unless
% a document genuinely needs another package.
\documentclass[11pt]{article}
\usepackage[a4paper,margin=2.7cm]{geometry}
\usepackage{amsmath,amssymb,amsthm,mathtools}
\usepackage[T1]{fontenc}
\usepackage{lmodern}
\usepackage{microtype}
\usepackage{xcolor}
\usepackage[colorlinks=true,linkcolor=blue!50!black,urlcolor=blue!50!black]{hyperref}
\allowdisplaybreaks

\newtheorem{lemma}{Lemma}
\newtheorem{proposition}{Proposition}
\theoremstyle{remark}
\newtheorem*{remark}{Remark}

\title{Cleo Bench Problem 3\\[1ex]\large How to find ${\large\int}_0^1\frac{\ln^3(1+x)\ln x}x\mathrm dx$}
\author{Derivation by Claude (Fable 5), closed-book\thanks{Problem originally
posed on Mathematics Stack Exchange
(\href{https://math.stackexchange.com/q/908108}{question 908108}, CC BY-SA),
famously answered by user Cleo. This derivation was produced independently,
offline, without access to the published answer, as part of the Cleo benchmark.}}
\date{July 2026}

\begin{document}
\maketitle

\section*{Solution: $\displaystyle\int_0^1\frac{\ln^3(1+x)\,\ln x}{x}\,dx$}

\section*{Problem}

Evaluate in closed form
\[
I=\int_0^1\frac{\ln^3(1+x)\,\ln x}{x}\,dx .
\]

\section*{Result}

\[
\boxed{\;I=\frac{99}{16}\,\zeta(5)+\frac{\pi^{2}}{2}\,\zeta(3)
-\frac{21}{4}\,\zeta(3)\ln^{2}2+\frac{\pi^{2}}{3}\,\ln^{3}2
-\frac{2}{5}\,\ln^{5}2-12\,\ln 2\,\operatorname{Li}_4\!\Big(\frac12\Big)
-12\,\operatorname{Li}_5\!\Big(\frac12\Big)\;}
\]

Equivalently, with $\zeta(2)=\pi^2/6$,
\[
I=\frac{99}{16}\zeta(5)+3\,\zeta(2)\zeta(3)-\frac{21}{4}\zeta(3)\ln^2 2
+2\,\zeta(2)\ln^3 2-\frac25\ln^5 2-12\ln 2\operatorname{Li}_4(\tfrac12)-12\operatorname{Li}_5(\tfrac12).
\]

Numerically $I=-0.05760941065877899523479148890918537945886000222364\ldots$

\section*{Derivation}

Throughout, $\eta(s)=\sum_{n\ge1}\frac{(-1)^{n-1}}{n^s}=(1-2^{1-s})\zeta(s)$, so
\[
\eta(2)=\tfrac12\zeta(2),\qquad \eta(3)=\tfrac34\zeta(3),\qquad
\eta(4)=\tfrac78\zeta(4),\qquad \eta(5)=\tfrac{15}{16}\zeta(5).
\]

For double series we write, for $s\ge1,t\ge1$,
\[
\begin{aligned}
&\zeta(s,t)=\sum_{n>m\ge1}\frac{1}{n^s m^t},\qquad
\zeta(\bar s,t)=\sum_{n>m\ge1}\frac{(-1)^n}{n^s m^t},\\
&\zeta(s,\bar t)=\sum_{n>m\ge1}\frac{(-1)^m}{n^s m^t},\qquad
\zeta(\bar s,\bar t)=\sum_{n>m\ge1}\frac{(-1)^{n+m}}{n^s m^t},
\end{aligned}
\]
and $\zeta(\bar s)=\sum_n(-1)^n/n^s=-\eta(s)$. Every double series above converges for the index
combinations used below: the inner sums are bounded (or $O(\log n)$), and the outer factor is
either $O(n^{-2})$ (absolute convergence after this bound) or alternating with terms of decreasing
modulus (Dirichlet's test). Rearrangements of double series are justified individually where they
occur (they are either between absolutely convergent arrangements or handled by Lemma 3.1).

We also use $H_n=\sum_{k\le n}\frac1k$, $H_n^{(r)}=\sum_{k\le n}\frac1{k^r}$.

\subsection*{0. Elementary lemmas}

\textbf{Lemma 0.1 (log moments).} For $n\ge1$, $a\ge0$, integrating the elementary antiderivative
$\int x^{n-1}\ln^a x\,dx=x^{n}\sum_{i=0}^{a}(-1)^{a-i}\frac{a!}{i!}\frac{\ln^i x}{n^{a+1-i}}$ (check by
differentiation) between the given limits:
\[
\int_0^1 x^{n-1}\ln^a x\,dx=\frac{(-1)^a a!}{n^{a+1}},\qquad
\int_0^{1/2}x^{n-1}\ln^a x\,dx=(-1)^a\,2^{-n}\sum_{i=0}^{a}\frac{a!}{i!}\frac{\ln^i 2}{n^{a+1-i}},
\]
\[
\int_{1/2}^{1}x^{n-1}\ln^a x\,dx=\frac{(-1)^a a!}{n^{a+1}}
-(-1)^a\,2^{-n}\sum_{i=0}^{a}\frac{a!}{i!}\frac{\ln^i 2}{n^{a+1-i}} .
\]
(All boundary terms at $x=0$ vanish since $x^n\ln^ix\to0$.)

\textbf{Lemma 0.2 (power series).} For $|x|<1$,
\[
-\ln(1-x)=\sum_{n\ge1}\frac{x^n}{n},\qquad
\ln^2(1-x)=2\sum_{n\ge2}\frac{H_{n-1}}{n}\,x^{n},\qquad
\ln^2(1+x)=2\sum_{n\ge2}(-1)^n\frac{H_{n-1}}{n}\,x^{n}.
\]
\textit{Proof.} Square the first series (Cauchy product, absolutely convergent for $|x|<1$):
the coefficient of $x^n$ is $\sum_{j+k=n}\frac{1}{jk}=\frac1n\sum_{j=1}^{n-1}\Big(\frac1j+\frac1{n-j}\Big)
=\frac{2H_{n-1}}{n}$, using $\frac{1}{j(n-j)}=\frac1n\big(\frac1j+\frac1{n-j}\big)$. Replace $x\to-x$
for the third identity. $\square$

\textbf{Lemma 0.3 (termwise integration).}
(a) \textit{Positive case.} If $f(x)=\sum c_nx^n$ on $[0,\tfrac12]$ with $c_n\ge0$ and $w\ge0$ is measurable,
then $\int_0^{1/2}wf=\sum c_n\int_0^{1/2}wx^n$ (monotone convergence).
(b) \textit{Alternating case.} Let $f(x)=\sum_{n\ge n_0}(-1)^n c_n x^n$ on $[0,1]$ with $c_n\downarrow 0$.
Then for every weight $w$ with $\int_0^1|w(x)|x^{n}dx<\infty$,
\[
\int_0^1 w f=\sum_{n}(-1)^nc_n\int_0^1 w(x)x^n dx,
\]
because the alternating-series remainder bound gives $|f(x)-S_N(x)|\le c_{N+1}x^{N+1}$ pointwise, so
$\big|\int w(f-S_N)\big|\le c_{N+1}\int_0^1 |w(x)|\,x^{N+1}dx\to0$ (for the weights used below,
$w=\ln^2x/x$ or $w=1/x$, Lemma 0.1 gives $\int_0^1|w|x^{N+1}dx = \tfrac{2}{(N+1)^3}$ or
$\tfrac{1}{N+1}$).
The coefficients $c_n=2H_{n-1}/n$ of $\ln^2(1+x)$ are decreasing for $n\ge2$ since
$H_{n-1}(n+1)-nH_n=H_{n-1}-1\ge0$. $\square$

\subsection*{1. Reduction of $I$ to half-interval integrals}

\textbf{Step 1.1 (integration by parts).} With $dv=\frac{\ln x}{x}dx$, $v=\frac{\ln^2x}{2}$, $u=\ln^3(1+x)$:
\begin{equation}
I=\Big[\tfrac12\ln^2x\,\ln^3(1+x)\Big]_0^1-\frac32\int_0^1\frac{\ln^2x\,\ln^2(1+x)}{1+x}\,dx
=-\frac32\,K,\qquad K:=\int_0^1\frac{\ln^2x\,\ln^2(1+x)}{1+x}\,dx, \tag{1}
\end{equation}
since $\ln^2 x\,\ln^3(1+x)\sim x^3\ln^2x\to0$ at $0$ and $\ln^21=0$ at $1$.
\textit{(Equation (1) will only be used as a consistency check; the evaluation below determines $I$
without it.)}

\textbf{Step 1.2 (transforming $K$).} Substitute $u=1+x$, then $u=1/v$ (so $v$ runs from $1$ to $\tfrac12$,
$\ln(u-1)=\ln\frac{1-v}{v}=\ln(1-v)-\ln v$, $\ln u=-\ln v$, $\frac{du}{u}=-\frac{dv}{v}$):
\[
K=\int_1^2\frac{\ln^2(u-1)\,\ln^2u}{u}\,du
=\int_{1/2}^{1}\frac{\big[\ln(1-v)-\ln v\big]^2\ln^2 v}{v}\,dv
=T_1-2T_2+\int_{1/2}^1\frac{\ln^4v}{v}dv,
\]
with
\[
T_1:=\int_{1/2}^{1}\frac{\ln^2(1-v)\,\ln^2v}{v}\,dv,\qquad
T_2:=\int_{1/2}^{1}\frac{\ln(1-v)\,\ln^3v}{v}\,dv,\qquad
\int_{1/2}^1\frac{\ln^4v}{v}dv=\frac{\ln^52}{5}.
\]
Hence
\begin{equation}
K=T_1-2T_2+\frac{\ln^52}{5}. \tag{2}
\end{equation}

\textbf{Step 1.3 ($T_2$ in closed form).} Expand $-\ln(1-v)=\sum_k v^k/k$ (positive terms); multiplying
by the nonnegative weight $-\ln^3v/v$ on $[\tfrac12,1]$, termwise integration is justified by
monotone convergence, and Lemma 0.1 gives
\[
T_2=-\sum_{k\ge1}\frac1k\int_{1/2}^1v^{k-1}\ln^3v\,dv
=\sum_{k\ge1}\frac{6}{k^5}-\sum_{k\ge1}\frac{2^{-k}}{k}\Big(\frac{\ln^32}{k}+\frac{3\ln^22}{k^2}
+\frac{6\ln2}{k^3}+\frac{6}{k^4}\Big),
\]
i.e.
\begin{equation}
T_2=6\zeta(5)-\ln^32\operatorname{Li}_2(\tfrac12)-3\ln^22\operatorname{Li}_3(\tfrac12)
-6\ln2\operatorname{Li}_4(\tfrac12)-6\operatorname{Li}_5(\tfrac12). \tag{3}
\end{equation}

\textbf{Step 1.4 ($T_1$ via the full interval).} By Lemma 0.2/0.3(a) and Lemma 0.1,
\[
\int_0^1\frac{\ln^2x\ln^2(1-x)}{x}\,dx
=2\sum_{n}\frac{H_{n-1}}{n}\cdot\frac{2}{n^3}=4\big(\textstyle\sum_n\frac{H_n}{n^4}-\zeta(5)\big),
\]
and Euler's evaluation $\sum_n H_n/n^4=3\zeta(5)-\zeta(2)\zeta(3)$ --- proved in Section 2 --- yields
\begin{equation}
\int_0^1\frac{\ln^2x\,\ln^2(1-x)}{x}\,dx=8\zeta(5)-4\zeta(2)\zeta(3). \tag{4}
\end{equation}
Therefore, writing $M:=\displaystyle\int_0^{1/2}\frac{\ln^2x\,\ln^2(1-x)}{x}\,dx$,
\[
T_1=\big(8\zeta(5)-4\zeta(2)\zeta(3)\big)-M .
\]
Combining with (2):
\begin{equation}
\boxed{\;K+M=8\zeta(5)-4\zeta(2)\zeta(3)-2T_2+\frac{\ln^52}{5}.\;}\tag{E1}
\end{equation}

\subsection*{2. Euler's non-alternating double zetas of weight $\le5$}

We need $\zeta(4,1)$ (for (4)) and, in Section 3, $\zeta(3,2),\zeta(2,3),\zeta(2,1)$ and
$\zeta(4)=\tfrac25\zeta(2)^2$. All proofs are by elementary rearrangements of \textbf{positive} double
series, so all rearrangements are legitimate.

\textbf{Lemma 2.1 (partial-fraction / Nielsen decomposition).} For integers $a,b\ge1$ and $j,k\ge1$,
\[
\frac{1}{j^ak^b}=\sum_{i=0}^{a-1}\binom{b-1+i}{i}\frac{1}{j^{a-i}\,(j+k)^{b+i}}
+\sum_{i=0}^{b-1}\binom{a-1+i}{i}\frac{1}{k^{b-i}\,(j+k)^{a+i}} .
\]
\textit{Proof.} Fix $n=j+k$ and decompose the rational function $r(j)=j^{-a}(n-j)^{-b}$ of $j$ into partial
fractions: it has poles of order $a$ at $j=0$ and order $b$ at $j=n$. The principal part at $j=0$ is
obtained from the Taylor expansion $(n-j)^{-b}=\sum_{i\ge0}\binom{b-1+i}{i}\frac{j^i}{n^{b+i}}$,
giving the coefficients of $j^{-(a-i)}$ as stated; symmetrically at $j=n$ with $k=n-j$. Since $r$
vanishes at $\infty$ faster than $1/j$, it equals the sum of its two principal parts. $\square$

Summing Lemma 2.1 over $j,k\ge1$ (positive terms), and using
$\sum_{j,k}\frac{1}{j^{c}(j+k)^{d}}=\sum_{n>j}\frac{1}{j^cn^d}=\zeta(d,c)$:
\begin{equation}
\zeta(a)\zeta(b)=\sum_{i=0}^{a-1}\binom{b-1+i}{i}\zeta(b+i,a-i)
+\sum_{i=0}^{b-1}\binom{a-1+i}{i}\zeta(a+i,b-i)\qquad(a,b\ge2). \tag{N}
\end{equation}

\textbf{Stuffle.} Splitting $\sum_{j,k}=\sum_{j>k}+\sum_{k>j}+\sum_{j=k}$ (positive terms),
\begin{equation}
\zeta(a)\zeta(b)=\zeta(a,b)+\zeta(b,a)+\zeta(a+b)\qquad(a,b\ge2). \tag{S}
\end{equation}

\textbf{Weight 4.} Consider $T:=\sum_{j,k\ge1}\frac{1}{j^2k(j+k)}$. Using
$\frac{1}{k(j+k)}=\frac1j\big(\frac1k-\frac1{j+k}\big)$ and $\sum_k\big(\frac1k-\frac1{j+k}\big)=H_j$:
\[
T=\sum_j\frac{H_j}{j^3}=\zeta(3,1)+\zeta(4).
\]
On the other hand $\frac{1}{jk}=\frac{1}{j+k}\big(\frac1j+\frac1k\big)$ applied inside $T$ gives
$T=\sum_{j,k}\frac{1}{j^2(j+k)^2}+\sum_{j,k}\frac{1}{jk(j+k)^2}$
and again $\frac{1}{jk}=\frac1{j+k}(\frac1j+\frac1k)$ in the second term:
$T=\zeta(2,2)+2\sum_{j,k}\frac{1}{j(j+k)^3}=\zeta(2,2)+2\zeta(3,1)$. Hence
\[
\zeta(2,2)=\zeta(4)-\zeta(3,1).
\]
Now (S) with $(2,2)$: $\zeta(2)^2=2\zeta(2,2)+\zeta(4)=3\zeta(4)-2\zeta(3,1)$, and (N) with $(2,2)$:
$\zeta(2)^2=2\big[\zeta(2,2)+2\zeta(3,1)\big]=2\zeta(4)+2\zeta(3,1)$. Subtracting,
\begin{equation}
\zeta(3,1)=\tfrac14\zeta(4),\qquad \zeta(2,2)=\tfrac34\zeta(4),\qquad
\zeta(2)^2=\tfrac52\zeta(4)\ \Longleftrightarrow\ \zeta(4)=\tfrac25\zeta(2)^2. \tag{5}
\end{equation}

\textbf{$\zeta(2,1)=\zeta(3)$.} Expanding $\ln^2(1-x)$ by Lemma 0.2 and integrating termwise
(positive terms, Lemma 0.1): $\int_0^1\frac{\ln^2(1-x)}{x}dx=2\sum_n\frac{H_{n-1}}{n^2}=2\zeta(2,1)$;
substituting $x\to1-x$, the same integral is $\int_0^1\frac{\ln^2x}{1-x}dx=\sum_n\frac{2}{n^3}=2\zeta(3)$.
Hence
\begin{equation}
\zeta(2,1)=\zeta(3). \tag{6}
\end{equation}

\textbf{Weight 5.} Consider $U:=\sum_{j,k\ge1}\frac{1}{j^2k^2(j+k)}$ (positive, convergent).
\textit{First evaluation.} Dividing the identity
$\frac{1}{j^2(j+k)}=\frac{1}{kj^2}-\frac{1}{k^2j}+\frac{1}{k^2(j+k)}$ (verify by clearing
denominators) by $k^2$:
\[
\frac{1}{j^2k^2(j+k)}=\frac{1}{j^2k^3}-\frac{1}{jk^4}+\frac{1}{k^4(j+k)} .
\]
Summing over $j,k$, the last two terms combine as
$\sum_k\frac1{k^4}\sum_j\big(\frac{1}{j+k}-\frac1j\big)=-\sum_k\frac{H_k}{k^4}
=-\zeta(4,1)-\zeta(5)$, so
\[
U=\zeta(2)\zeta(3)-\zeta(4,1)-\zeta(5).
\]
\textit{Second evaluation.} $\frac{1}{jk}=\frac1{j+k}(\frac1j+\frac1k)$ twice:
$U=2\sum_{j,k}\frac{1}{jk^2(j+k)^2}=2\Big[\sum_{j,k}\frac{1}{k^2(j+k)^3}
+\sum_{j,k}\frac{1}{jk(j+k)^3}\Big]
=2\zeta(3,2)+4\sum_{j,k}\frac{1}{j(j+k)^4}=2\zeta(3,2)+4\zeta(4,1)$.
Equating:
\begin{equation}
2\zeta(3,2)+5\zeta(4,1)=\zeta(2)\zeta(3)-\zeta(5). \tag{R1}
\end{equation}
The stuffle (S) and Nielsen (N) instances for $(a,b)=(2,3)$ read
\begin{equation}
\zeta(2,3)+\zeta(3,2)+\zeta(5)=\zeta(2)\zeta(3), \tag{R2}
\end{equation}
\begin{equation}
\zeta(2,3)+3\zeta(3,2)+6\zeta(4,1)=\zeta(2)\zeta(3). \tag{R3}
\end{equation}
Subtracting (R2) from (R3): $2\zeta(3,2)+6\zeta(4,1)=\zeta(5)$; subtracting (R1) from this:
\begin{equation}
\zeta(4,1)=2\zeta(5)-\zeta(2)\zeta(3),\qquad
\zeta(3,2)=3\zeta(2)\zeta(3)-\tfrac{11}{2}\zeta(5),\qquad
\zeta(2,3)=\tfrac92\zeta(5)-2\zeta(2)\zeta(3). \tag{7}
\end{equation}
In particular $\sum_n\frac{H_n}{n^4}=\zeta(4,1)+\zeta(5)=3\zeta(5)-\zeta(2)\zeta(3)$, which proves (4).

\subsection*{3. Alternating double zetas: $\zeta(\bar2,1)$ and $\zeta(\bar4,1)$}

Here the three sign patterns $(\bar s,t),(s,\bar t),(\bar s,\bar t)$ are pinned by three families of
identities.

\textbf{(a) Signed stuffles.} Exactly as (S), splitting the product of two convergent single series into
$\{j>k\},\{k>j\},\{j=k\}$: if $\alpha,\beta\in\{1,-1\}$ and $(\text{outer exponent},\text{sign})$
makes each series converge, then with $\zeta^{\alpha}(a):=\sum_j\alpha^j/j^a$
($\zeta^{-1}(a)=-\eta(a)$, $\zeta^{+1}(a)=\zeta(a)$),
\[
\zeta^{\alpha}(a)\,\zeta^{\beta}(b)
=\sum_{j>k}\frac{\alpha^j\beta^k}{j^ak^b}+\sum_{k>j}\frac{\beta^k\alpha^j}{k^bj^a}
+\sum_{j}\frac{(\alpha\beta)^j}{j^{a+b}}.
\]
\textit{Justification:} multiply the partial sums up to $N$ (an exact identity), and let $N\to\infty$; each
of the three partial double sums converges to the corresponding double zeta (outer index bounded by
$N$; the inner sums are bounded, and the outer series converge absolutely or by Dirichlet's test).

\textbf{(b) Signed Nielsen relations.} Multiply Lemma 2.1 by $\sigma^j\tau^k$ ($\sigma,\tau\in\{\pm1\}$)
and sum over $j,k\ge1$ as an iterated sum ($k$ outer, $j$ inner). On the left, the inner sum over
$j$ is the fixed convergent series $\sum_j\sigma^j j^{-a}$ (alternating when $\sigma=-1$, $a=1$),
so the iterated sum equals the product $\zeta^{\sigma}(a)\zeta^{\tau}(b)$ of the two single series.
On the right, each term is regrouped via $n=j+k$ by the following lemma.

\textbf{Lemma 3.1 (regrouping).} Let $\sigma,\tau\in\{\pm1\}$, $c\ge1$, $d\ge1$, $c+d\ge3$, and exclude
$(d,\sigma)=(1,+1)$. Then
\[
\sum_{k\ge1}\frac{\tau^k}{k^{c}}\sum_{j\ge1}\frac{\sigma^j}{(j+k)^{d}}
=\sum_{n>k\ge1}\frac{\sigma^{n}(\sigma\tau)^{k}}{n^{d}\,k^{c}},
\]
both sides converging.
\textit{Proof.} For fixed $k$, reindex $n=j+k$ (a bijective, order-preserving reindexing of a single
series, always valid): $\sum_j\sigma^j(j+k)^{-d}=\sigma^{-k}\sum_{n>k}\sigma^n n^{-d}$. So the left
side is $\sum_k(\sigma\tau)^k k^{-c}\sum_{n>k}\sigma^n n^{-d}$. If $d\ge2$ everything is absolutely
convergent ($\sum_{n>k}n^{-d}\le Ck^{1-d}$, and $\sum_k k^{-c}\,k^{1-d}<\infty$ since $c+d\ge3$), and
swapping to $\sum_n\sum_{k<n}$ is Fubini. If $d=1$, $\sigma=-1$: let
$P_K:=\sum_{k\le K}(\sigma\tau)^kk^{-c}\sum_{n>k}\sigma^nn^{-1}$. Splitting the inner tail at $n=K$,
\[
P_K=\sum_{n\le K}\frac{\sigma^n}{n}\sum_{k<n}\frac{(\sigma\tau)^k}{k^c}
+\Big(\sum_{k\le K}\frac{(\sigma\tau)^k}{k^c}\Big)\sum_{n>K}\frac{\sigma^n}{n},
\]
(the first term is a finite triangular swap). The last factor is an alternating tail,
$\big|\sum_{n>K}\sigma^n/n\big|\le\frac1{K+1}$, while $\big|\sum_{k\le K}(\sigma\tau)^kk^{-c}\big|\le
1+\log K$; the product tends to $0$. In the first term note that $d=1$ together with $c+d\ge3$
forces $c\ge2$, so the inner sums $\sum_{k<n}(\sigma\tau)^kk^{-c}$ are bounded and converge; by
Dirichlet's test (alternating factor $\sigma^n/n$ with $\sigma=-1$, bounded inner sums) the first
term converges as $K\to\infty$ to the right side. $\square$

The right side of Lemma 3.1 is the double zeta with outer sign $\sigma$ and inner sign
$\sigma\tau$ attached to exponents $(d,c)$. Applying it to each Nielsen term
$1/(j^{a-i}(j+k)^{b+i})$ (with the roles of $j,k$ as written) yields, for $\sigma,\tau$ not both
$+1$,
\[
\zeta^{\sigma}(a)\,\zeta^{\tau}(b)=\sum_{i=0}^{a-1}\binom{b-1+i}{i}\,
\zeta^{\tau,\sigma\tau}(b+i,\,a-i)+\sum_{i=0}^{b-1}\binom{a-1+i}{i}\,
\zeta^{\sigma,\sigma\tau}(a+i,\,b-i),
\]
where $\zeta^{\alpha,\beta}(s,t):=\sum_{n>m}\alpha^n\beta^m n^{-s}m^{-t}$, valid whenever every
term is covered by Lemma 3.1 --- true in all instances used below (each term has weight $3$ or $5$
and outer exponent $\ge2$, or outer exponent $1$ with outer sign $-1$).

\textbf{(c) Doubling (distribution).} For $s\ge2$, $t\ge1$, summing the four sign patterns termwise
(a finite linear combination of convergent series) gives
\begin{equation}
\zeta(s,t)+\zeta(\bar s,t)+\zeta(s,\bar t)+\zeta(\bar s,\bar t)
=\sum_{n>m}\frac{(1+(-1)^n)(1+(-1)^m)}{n^sm^t}
=4\!\!\sum_{\substack{n>m\\ n,m\ \mathrm{even}}}\!\!\frac{1}{n^sm^t}
=2^{2-s-t}\,\zeta(s,t). \tag{D}
\end{equation}

\textbf{Weight 3.} The five unknowns $\zeta(\bar2,1),\zeta(2,\bar1),\zeta(\bar2,\bar1),\zeta(\bar1,2),
\zeta(\bar1,\bar2)$ satisfy:

\begin{center}
\begin{tabular}{|c|l|l|}
\hline
\# & identity & source \\
\hline
1 & $\zeta(\bar1,2)+\zeta(2,\bar1)-\eta(3)=-\ln2\,\zeta(2)$ & stuffle $\zeta^{-}(1)\zeta^{+}(2)$ \\
2 & $\zeta(\bar1,\bar2)+\zeta(\bar2,\bar1)+\zeta(3)=\ln2\,\eta(2)$ & stuffle $\zeta^{-}(1)\zeta^{-}(2)$ \\
3 & $\zeta(2,\bar1)+\zeta(\bar1,\bar2)+\zeta(\bar2,\bar1)=-\ln2\,\zeta(2)$ & Nielsen $(a,b)=(1,2)$, $(\sigma,\tau)=(-,+)$ \\
4 & $2\zeta(\bar2,1)+\zeta(\bar1,2)=\ln2\,\eta(2)$ & Nielsen $(1,2)$, $(\sigma,\tau)=(-,-)$ \\
5 & $\zeta(\bar2,1)+\zeta(2,\bar1)+\zeta(\bar2,\bar1)=-\tfrac12\zeta(2,1)=-\tfrac12\zeta(3)$ & doubling (D) + (6) \\
\hline
\end{tabular}
\end{center}

(For row 3: $(a,b;\sigma,\tau)=(1,2;-,+)$ in (b) gives left side $\big(\sum_j\tfrac{(-1)^j}{j}\big)\zeta(2)$
and right side $\zeta^{\tau,\sigma\tau}(2,1)+\zeta^{\sigma,\sigma\tau}(1,2)+\zeta^{\sigma,\sigma\tau}(2,1)
=\zeta(2,\bar1)+\zeta(\bar1,\bar2)+\zeta(\bar2,\bar1)$; row 4 similarly with $(\sigma,\tau)=(-,-)$.)

This linear system has full rank $5$ (a direct check of the $5\times5$ coefficient matrix), and its
unique solution is
\[
\zeta(\bar2,1)=\tfrac18\zeta(3),\qquad
\zeta(2,\bar1)=\zeta(3)-\tfrac32\ln2\,\zeta(2),\qquad
\zeta(\bar2,\bar1)=-\tfrac{13}{8}\zeta(3)+\tfrac32\ln2\,\zeta(2),
\]
\begin{equation}
\zeta(\bar1,2)=-\tfrac14\zeta(3)+\tfrac12\ln2\,\zeta(2),\qquad
\zeta(\bar1,\bar2)=\tfrac58\zeta(3)-\ln2\,\zeta(2). \tag{8}
\end{equation}
(Substituting these values back into rows 1--5 verifies the solution; e.g. row 5:
$\tfrac18+1-\tfrac{13}8=-\tfrac12$ for the $\zeta(3)$ coefficients and $-\tfrac32+\tfrac32=0$ for
$\ln2\,\zeta(2)$.) In particular
\begin{equation}
J:=\int_0^1\frac{\ln^2(1+u)}{u}\,du
=2\sum_{n}\frac{(-1)^nH_{n-1}}{n^2}=2\,\zeta(\bar2,1)=\frac{\zeta(3)}{4}, \tag{9}
\end{equation}
where the termwise integration uses Lemma 0.3(b) with weight $1/u$, and
$\sum_n(-1)^nH_{n-1}/n^2=\sum_{n>m}(-1)^n/(n^2m)=\zeta(\bar2,1)$.

\textbf{Weight 5.} The eleven unknowns are the signed versions of $(4,1),(3,2),(2,3)$ (three sign
patterns each) and $(\bar1,4),(\bar1,\bar4)$. The same three mechanisms give thirteen identities
(here $\zeta(\bar5)=-\eta(5)$; right-hand sides use $\eta$-values and (7)):

\begin{center}
\begin{tabular}{|c|l|l|}
\hline
\# & identity & source \\
\hline
1 & $\zeta(\bar1,4)+\zeta(4,\bar1)-\eta(5)=-\ln2\,\zeta(4)$ & st. $\zeta^{-}(1)\zeta^{+}(4)$ \\
2 & $\zeta(\bar1,\bar4)+\zeta(\bar4,\bar1)+\zeta(5)=\ln2\,\eta(4)$ & st. $\zeta^{-}(1)\zeta^{-}(4)$ \\
3 & $\zeta(\bar2,3)+\zeta(3,\bar2)-\eta(5)=-\eta(2)\zeta(3)$ & st. $\zeta^{-}(2)\zeta^{+}(3)$ \\
4 & $\zeta(2,\bar3)+\zeta(\bar3,2)-\eta(5)=-\zeta(2)\eta(3)$ & st. $\zeta^{+}(2)\zeta^{-}(3)$ \\
5 & $\zeta(\bar2,\bar3)+\zeta(\bar3,\bar2)+\zeta(5)=\eta(2)\eta(3)$ & st. $\zeta^{-}(2)\zeta^{-}(3)$ \\
6 & $\zeta(4,\bar1)+\zeta(\bar1,\bar4)+\zeta(\bar2,\bar3)+\zeta(\bar3,\bar2)+\zeta(\bar4,\bar1)=-\ln2\,\zeta(4)$ & Ni. $(1,4;-,+)$ \\
7 & $2\zeta(\bar4,1)+\zeta(\bar1,4)+\zeta(\bar2,3)+\zeta(\bar3,2)=\ln2\,\eta(4)$ & Ni. $(1,4;-,-)$ \\
8 & $\zeta(3,\bar2)+3\zeta(4,\bar1)+\zeta(\bar2,\bar3)+2\zeta(\bar3,\bar2)+3\zeta(\bar4,\bar1)=-\eta(2)\zeta(3)$ & Ni. $(2,3;-,+)$ \\
9 & $\zeta(\bar3,\bar2)+3\zeta(\bar4,\bar1)+\zeta(2,\bar3)+2\zeta(3,\bar2)+3\zeta(4,\bar1)=-\zeta(2)\eta(3)$ & Ni. $(2,3;+,-)$ \\
10 & $\zeta(\bar2,3)+3\zeta(\bar3,2)+6\zeta(\bar4,1)=\eta(2)\eta(3)$ & Ni. $(2,3;-,-)$ \\
11 & $\zeta(\bar4,1)+\zeta(4,\bar1)+\zeta(\bar4,\bar1)=-\tfrac78\zeta(4,1)$ & (D) \\
12 & $\zeta(\bar3,2)+\zeta(3,\bar2)+\zeta(\bar3,\bar2)=-\tfrac78\zeta(3,2)$ & (D) \\
13 & $\zeta(\bar2,3)+\zeta(2,\bar3)+\zeta(\bar2,\bar3)=-\tfrac78\zeta(2,3)$ & (D) \\
\hline
\end{tabular}
\end{center}

(In rows 6--10, the general recipe of (b) is used: for $(a,b;\sigma,\tau)$ the right side of Lemma 2.1
contributes $\sum_{i=0}^{a-1}\binom{b-1+i}{i}\,\zeta^{\tau,\sigma\tau}(b+i,a-i)
+\sum_{i=0}^{b-1}\binom{a-1+i}{i}\,\zeta^{\sigma,\sigma\tau}(a+i,b-i)$, where
$\zeta^{\alpha,\beta}(s,t)=\sum_{n>m}\alpha^n\beta^m n^{-s}m^{-t}$.)

This $13\times11$ system has rank $11$ (checked by Gaussian elimination over $\mathbb{Q}$), hence a
unique solution, which one verifies by substitution:
\begin{equation}
\begin{aligned}
&\zeta(\bar4,1)=-\tfrac{29}{32}\zeta(5)+\tfrac12\zeta(2)\zeta(3), &
&\zeta(4,\bar1)=2\zeta(5)-\tfrac38\zeta(2)\zeta(3)-\tfrac{15}{8}\zeta(4)\ln2,\\
&\zeta(\bar4,\bar1)=-\tfrac{91}{32}\zeta(5)+\tfrac34\zeta(2)\zeta(3)+\tfrac{15}{8}\zeta(4)\ln2, &
&\zeta(\bar3,2)=\tfrac{41}{32}\zeta(5)-\tfrac58\zeta(2)\zeta(3),\\
&\zeta(3,\bar2)=-\tfrac{21}{32}\zeta(5)+\tfrac14\zeta(2)\zeta(3), &
&\zeta(\bar3,\bar2)=\tfrac{67}{16}\zeta(5)-\tfrac94\zeta(2)\zeta(3),\\
&\zeta(\bar2,3)=\tfrac{51}{32}\zeta(5)-\tfrac34\zeta(2)\zeta(3), &
&\zeta(2,\bar3)=-\tfrac{11}{32}\zeta(5)-\tfrac18\zeta(2)\zeta(3),\\
&\zeta(\bar2,\bar3)=-\tfrac{83}{16}\zeta(5)+\tfrac{21}{8}\zeta(2)\zeta(3), &
&\zeta(\bar1,4)=-\tfrac{17}{16}\zeta(5)+\tfrac38\zeta(2)\zeta(3)+\tfrac78\zeta(4)\ln2,\\
&\zeta(\bar1,\bar4)=\tfrac{59}{32}\zeta(5)-\tfrac34\zeta(2)\zeta(3)-\zeta(4)\ln2,&&
\end{aligned}\tag{10}
\end{equation}
using $\zeta(2)^2=\tfrac52\zeta(4)$ from (5) to express the $\ln2\,\zeta(2)^2$ terms via $\zeta(4)\ln2$.
The value we need is the first one. Since
\[
A:=\sum_{n\ge1}\frac{(-1)^{n-1}H_n}{n^4}
=\eta(5)+\sum_{n>m}\frac{(-1)^{n-1}}{n^4m}=\eta(5)-\zeta(\bar4,1),
\]
we obtain the classical
\begin{equation}
A=\frac{15}{16}\zeta(5)+\frac{29}{32}\zeta(5)-\frac12\zeta(2)\zeta(3)
=\frac{59}{32}\,\zeta(5)-\frac{1}{2}\,\zeta(2)\zeta(3). \tag{11}
\end{equation}
\textit{(Remark: the same table gives
$\sum(-1)^{n-1}H^{(2)}_n/n^3=\eta(5)-\zeta(\bar3,2)=\tfrac58\zeta(2)\zeta(3)-\tfrac{11}{32}\zeta(5)$,
a known cross-check.)}

\subsection*{4. $\operatorname{Li}_2(\tfrac12)$ and $\operatorname{Li}_3(\tfrac12)$}

\textbf{Lemma 4.1.} $\operatorname{Li}_2(x)+\operatorname{Li}_2(1-x)=\zeta(2)-\ln x\ln(1-x)$ on $(0,1)$:
both sides have derivative $-\frac{\ln(1-x)}{x}+\frac{\ln x}{1-x}$ and limit $\zeta(2)$ as
$x\to0^+$. At $x=\tfrac12$:
\begin{equation}
\operatorname{Li}_2(\tfrac12)=\tfrac12\zeta(2)-\tfrac12\ln^22. \tag{12}
\end{equation}

\textbf{Lemma 4.2.} Let $\tau:=\sum_{n\ge2}\frac{H_{n-1}}{2^n n^2}$. Then
\[
\tau=\frac{\zeta(3)}{8}-\frac{\ln^32}{6}
\qquad\text{and}\qquad
\tau=\zeta(3)-\frac{\ln^32}{2}-\ln2\operatorname{Li}_2(\tfrac12)-\operatorname{Li}_3(\tfrac12),
\]
and consequently
\begin{equation}
\operatorname{Li}_3(\tfrac12)=\tfrac78\zeta(3)-\tfrac12\zeta(2)\ln2+\tfrac16\ln^32. \tag{13}
\end{equation}
\textit{Proof.} By Lemma 0.2/0.3(a), $\tau=\int_0^{1/2}\frac{\ln^2(1-t)}{2t}\,dt$.
(i) Substitute $t=\frac{u}{1+u}$ (increasing bijection $[0,1]\to[0,\tfrac12]$), so $1-t=\frac1{1+u}$,
$\frac{dt}{t}=\frac{du}{u(1+u)}=\frac{du}{u}-\frac{du}{1+u}$:
\[
\tau=\frac12\int_0^1\ln^2(1+u)\Big(\frac1u-\frac1{1+u}\Big)du
=\frac12 J-\frac12\cdot\frac{\ln^32}{3}=\frac{\zeta(3)}{8}-\frac{\ln^32}{6},
\]
by (9). (ii) Alternatively split $\int_0^{1/2}=\int_0^1-\int_{1/2}^1$: the full integral is
$\int_0^1\frac{\ln^2(1-t)}{2t}dt=\zeta(2,1)=\zeta(3)$ by (6), while in
$\int_{1/2}^1$ substitute $t\to1-t$ and integrate termwise (Lemma 0.1, positive terms):
\[
\int_{1/2}^{1}\frac{\ln^2(1-t)}{2t}\,dt=\int_0^{1/2}\frac{\ln^2t}{2(1-t)}\,dt
=\frac12\sum_{n\ge1}2^{-n}\Big(\frac{\ln^22}{n}+\frac{2\ln2}{n^2}+\frac{2}{n^3}\Big)
=\frac{\ln^32}2+\ln2\operatorname{Li}_2(\tfrac12)+\operatorname{Li}_3(\tfrac12).
\]
Equating (i) and (ii) and inserting (12) gives (13). $\square$

\subsection*{5. The Landen equation for $M$}

Recall $M=\displaystyle\int_0^{1/2}\frac{\ln^2t\,\ln^2(1-t)}{t}\,dt$ from Step 1.4. Apply the same
substitution $t=\frac{u}{1+u}$ as in Lemma 4.2: $\ln(1-t)=-\ln(1+u)$, $\ln t=\ln u-\ln(1+u)$,
$\frac{dt}{t}=\frac{du}{u}-\frac{du}{1+u}$. Hence
\[
M=\int_0^1\Big[\ln^2u\,\ln^2(1+u)-2\ln u\,\ln^3(1+u)+\ln^4(1+u)\Big]
\Big(\frac1u-\frac1{1+u}\Big)du .
\]
All six resulting integrals converge separately; name them:
\begin{equation}
M=\big(\mathcal I-2X_1+X_2\big)-\big(Y_0-2Y_1+Y_2\big), \tag{14}
\end{equation}
\[
\mathcal I=\int_0^1\frac{\ln^2u\ln^2(1+u)}{u}du,\quad
X_1=\int_0^1\frac{\ln u\ln^3(1+u)}{u}du,\quad
X_2=\int_0^1\frac{\ln^4(1+u)}{u}du,
\]
\[
Y_0=\int_0^1\frac{\ln^2u\ln^2(1+u)}{1+u}du,\quad
Y_1=\int_0^1\frac{\ln u\ln^3(1+u)}{1+u}du,\quad
Y_2=\int_0^1\frac{\ln^4(1+u)}{1+u}du .
\]

Now evaluate or identify each piece.

\begin{itemize}
\item \textbf{$X_1=I$} --- it \textit{is} the integral we are computing (same integrand, $u\leftrightarrow x$).
\item \textbf{$Y_0=K$} --- the integral from Step 1.1, by definition.
\item \textbf{$Y_2=\dfrac{\ln^52}{5}$} --- immediate.
\item \textbf{$Y_1=-\dfrac{X_2}{4}$}: integrate by parts with $d\big[\tfrac14\ln^4(1+u)\big]
  =\frac{\ln^3(1+u)}{1+u}du$:
  \[
  Y_1=\Big[\frac{\ln^4(1+u)}{4}\ln u\Big]_0^1-\frac14\int_0^1\frac{\ln^4(1+u)}{u}\,du=-\frac{X_2}4 ,
  \]
  the boundary vanishing at both ends ($\ln 1=0$; $\ln^4(1+u)\ln u\sim u^4\ln u\to0$).
\item \textbf{$\mathcal I=4\big(\eta(5)-A\big)$}: by Lemma 0.2 and Lemma 0.3(b) (weight $\ln^2u/u$),
  \[
  \mathcal I=2\sum_{n\ge2}\frac{(-1)^nH_{n-1}}{n}\int_0^1u^{n-1}\ln^2u\,du
  =4\sum_{n}\frac{(-1)^nH_{n-1}}{n^4}
  =4\Big(\eta(5)-A\Big),
  \]
  since $\sum(-1)^nH_{n-1}/n^4=\sum(-1)^nH_n/n^4-\sum(-1)^n/n^5=-A+\eta(5)$. With (11),
  \begin{equation}
  \mathcal I=4\eta(5)-4A=\frac{15}{4}\zeta(5)-\frac{59}{8}\zeta(5)+2\zeta(2)\zeta(3)
  =2\zeta(2)\zeta(3)-\frac{29}{8}\zeta(5). \tag{15}
  \end{equation}
\item \textbf{$X_2$ explicitly}: substitute $v=1+u$, then $v=1/w$ (as in Step 1.2):
  \[
  X_2=\int_1^2\frac{\ln^4v}{v-1}\,dv=\int_{1/2}^1\frac{\ln^4w}{w(1-w)}\,dw
  =\int_{1/2}^1\frac{\ln^4w}{w}dw+\int_{1/2}^1\frac{\ln^4w}{1-w}\,dw .
  \]
  The first integral is $\tfrac15\ln^52$. For the second, $\int_0^1\frac{\ln^4w}{1-w}dw
  =\sum_n\frac{24}{n^5}=24\zeta(5)$ and, termwise by Lemma 0.1 (positive terms),
  \[
  \int_0^{1/2}\frac{\ln^4w}{1-w}\,dw
  =\sum_{n\ge1}2^{-n}\Big(\frac{\ln^42}{n}+\frac{4\ln^32}{n^2}+\frac{12\ln^22}{n^3}
  +\frac{24\ln2}{n^4}+\frac{24}{n^5}\Big)
  \]
  \[
  =\ln^52+4\ln^32\operatorname{Li}_2(\tfrac12)+12\ln^22\operatorname{Li}_3(\tfrac12)
  +24\ln2\operatorname{Li}_4(\tfrac12)+24\operatorname{Li}_5(\tfrac12).
  \]
  Hence
  \begin{equation}
  X_2=24\zeta(5)-\tfrac45\ln^52-4\ln^32\operatorname{Li}_2(\tfrac12)
  -12\ln^22\operatorname{Li}_3(\tfrac12)-24\ln2\operatorname{Li}_4(\tfrac12)
  -24\operatorname{Li}_5(\tfrac12). \tag{16}
  \end{equation}
  (Comparing with (3), $X_2=4T_2-\tfrac45\ln^52$, a useful cross-check.)
\end{itemize}

Substituting all pieces into (14):
\begin{equation}
\boxed{\;M+K+2I=4\big(\eta(5)-A\big)+\frac{X_2}{2}-\frac{\ln^52}{5}.\;}\tag{E2}
\end{equation}

\subsection*{6. Solving for $I$}

Subtract (E1) from (E2); $M$ and $K$ cancel:
\[
2I=4\big(\eta(5)-A\big)+\frac{X_2}{2}-\frac{\ln^52}{5}
-8\zeta(5)+4\zeta(2)\zeta(3)+2T_2-\frac{\ln^52}{5},
\]
\begin{equation}
I=2\big(\eta(5)-A\big)+\frac{X_2}{4}+T_2-4\zeta(5)+2\zeta(2)\zeta(3)-\frac{\ln^52}{5}. \tag{17}
\end{equation}
Insert (11), (3), (16):
\[
2\big(\eta(5)-A\big)=\tfrac{15}{8}\zeta(5)-\tfrac{59}{16}\zeta(5)+\zeta(2)\zeta(3)
=-\tfrac{29}{16}\zeta(5)+\zeta(2)\zeta(3),
\]
\[
\frac{X_2}{4}+T_2=12\zeta(5)-\tfrac15\ln^52-2\ln^32\operatorname{Li}_2(\tfrac12)
-6\ln^22\operatorname{Li}_3(\tfrac12)-12\ln2\operatorname{Li}_4(\tfrac12)
-12\operatorname{Li}_5(\tfrac12).
\]
Therefore
\[
I=\Big(-\tfrac{29}{16}-4+12\Big)\zeta(5)+3\zeta(2)\zeta(3)-\tfrac25\ln^52
-2\ln^32\operatorname{Li}_2(\tfrac12)-6\ln^22\operatorname{Li}_3(\tfrac12)
-12\ln2\operatorname{Li}_4(\tfrac12)-12\operatorname{Li}_5(\tfrac12),
\]
and $-\tfrac{29}{16}+8=\tfrac{99}{16}$. Finally, expanding with (12) and (13),
\[
-2\ln^32\operatorname{Li}_2(\tfrac12)=-\zeta(2)\ln^32+\ln^52,\qquad
-6\ln^22\operatorname{Li}_3(\tfrac12)=-\tfrac{21}{4}\zeta(3)\ln^22+3\zeta(2)\ln^32-\ln^52,
\]
so the $\ln^52$ corrections cancel and the $\zeta(2)\ln^32$ terms add to $+2\zeta(2)\ln^32$:
\[
I=\frac{99}{16}\zeta(5)+3\zeta(2)\zeta(3)-\frac{21}{4}\zeta(3)\ln^22
+2\zeta(2)\ln^32-\frac25\ln^52
-12\ln2\operatorname{Li}_4(\tfrac12)-12\operatorname{Li}_5(\tfrac12),
\]
which is the boxed Result (using $\zeta(2)=\pi^2/6$).

\textbf{Consistency checks.} (a) Equation (1), $I=-\tfrac32K$, was never used above; direct quadrature
of both integrals confirms it to 110 digits,
$K=0.0384062737725193301565276592727902529725733\ldots=-\tfrac23 I$, so the overdetermined system
(1), (E1), (E2) is consistent. (b) The value
$M=2\ln^22\,\tau+4\ln2\sum_n\frac{H_{n-1}}{2^nn^3}+4\sum_n\frac{H_{n-1}}{2^nn^4}$ obtained from its
defining series (Lemma 0.1) matches direct quadrature of $M$ to more than 100 digits.

\section*{Numerical verification}

All computations with mpmath at 110 significant digits.

Direct numerical integration (\texttt{mp.quad} with splits at $\tfrac14,\tfrac12,\tfrac34$):
\begin{verbatim}
I_quad = -0.05760941065877899523479148890918537945886000222364830713296
          23962670128085772005803934336413232954600867954...
\end{verbatim}
Closed form $\frac{99}{16}\zeta(5)+3\zeta(2)\zeta(3)-\frac{21}{4}\zeta(3)\ln^22+2\zeta(2)\ln^32
-\frac25\ln^52-12\ln2\operatorname{Li}_4(\frac12)-12\operatorname{Li}_5(\frac12)$:
\begin{verbatim}
I_cf   = -0.05760941065877899523479148890918537945886000222364830713296
          23962670128085772005803934336413232954600867954...
\end{verbatim}
Absolute difference $1.11\cdot10^{-110}$ (relative $1.9\cdot10^{-109}$), i.e. \textbf{agreement to
$\approx108$ significant digits} (far beyond the required 25).

Independent second check: the nested-sum representation
$I=3\sum_{n\ge1}\frac{(-1)^n\big(H_{n-1}^2-H_{n-1}^{(2)}\big)}{n^3}$
(obtained by expanding $\ln^3(1+x)$; equivalent to $I=6\,\zeta(\bar3,1,1)$), summed with mpmath's
alternating-series acceleration (\texttt{mp.nsum(..., method='a')}), agrees with \texttt{I\_quad} to
$1.3\cdot10^{-63}$ at 60-digit working precision.

Every displayed intermediate identity --- (3), (4), (5)--(13), (15), (16), (E1), (E2) --- was verified
numerically to at least 55 significant digits, and the linear systems of Sections 2--3 were solved
exactly (rational arithmetic; ranks $3$, $5$ and $11$ as claimed, solutions verified by
substitution).

As an additional independent confirmation, the full algebra of alternating multiple zeta values of
weight $\le5$ was solved by computer over $\mathbb{Q}$ (relations: series stuffles, iterated-integral
shuffles, the doubling/distribution relations, and the duality/Landen substitutions --- each family
numerically validated to $10^{-55}$ before use), and it reproduces exactly the same closed form for
$I=6\zeta(\bar3,1,1)$, as well as the table (10) and the auxiliary constants
$\sum_n H_n/(2^nn^3)=\operatorname{Li}_4(\tfrac12)+\tfrac{\pi^4}{720}-\tfrac18\zeta(3)\ln2+\tfrac1{24}\ln^42$,
$\sum_n H_n/(2^nn^4)=2\operatorname{Li}_5(\tfrac12)+\ln2\operatorname{Li}_4(\tfrac12)
+\tfrac1{32}\zeta(5)-\tfrac{\pi^2}{12}\zeta(3)+\tfrac12\zeta(3)\ln^22-\tfrac{\pi^4}{720}\ln2
-\tfrac{\pi^2}{36}\ln^32+\tfrac1{40}\ln^52$ (not needed above, but classical).

\section*{Notes}

\begin{itemize}
\item The derivation is self-contained: the only external inputs are Euler's values
  $\zeta(2)=\pi^2/6$ (used only for cosmetic rewriting) and the elementary convergence facts
  quoted; $\zeta(4)=\tfrac25\zeta(2)^2$ is derived in (5).
\item Rearrangement/termwise-integration steps are justified where they occur: positive-term series use
  monotone convergence; the two alternating expansions use the remainder bound of Lemma 0.3(b);
  the signed stuffle/Nielsen identities are proved through partial sums with absolutely convergent
  regrouped pieces (weights are high enough that every regrouped double series converges
  absolutely, except the explicitly-handled outer alternating sums).
\item The linear systems in Sections 2 and 3 are small and over $\mathbb{Q}$; their unique solvability
  is a finite rational-arithmetic fact (ranks verified), and the displayed solutions can be checked
  by direct substitution into each listed identity.
\item The key structural trick is Step 5: the Landen map $t\mapsto t/(1-t)$ (equivalently
  $t=u/(1+u)$) sends the half-interval integral $M$ to a combination of $[0,1]$-integrals in which
  the unknown $I$ itself reappears ($X_1=I$), while the companion equation (E1) contains the same
  pair $M+K$; subtracting eliminates both auxiliary integrals. The polylogarithms
  $\operatorname{Li}_4(\tfrac12),\operatorname{Li}_5(\tfrac12)$ enter only through the completely
  elementary expansions (3) and (16), and are not further reducible (they are standard basis
  constants at weight 4 and 5, level 2).
\end{itemize}

\end{document}
